% !PRINTED 646
\weberpage{646}
If $t$ denotes a new variable, then
\[
N(\delta t - \delta_0) = D E(x t - x_0,\ y t - y_0, \ldots),
\]
and hence, if we arrange in descending powers of $t$ and cancel the factor $N(\delta) = D E$ on both sides:
\weq{4}{N\left(t - \frac{\delta_0}{\delta}\right) = t^n + \overline{C}_1 t^{n-1} + \overline{C}_2 t^{n-2} + \cdots,}
in which $\overline{C}_1, \overline{C}_2 \ldots$ have no denominator other than $E$, and consequently are integral functionals in $\overline{Z}$. Accordingly, by the definition of \S\,179, 6., $\delta_0 : \delta$ is likewise an integral functional, as was to be proved.

Accordingly we obtain the following definitions:

1. \gesperrt{The functional $\delta = \alpha x + \beta y + \cdots$, and every functional associated with $\delta$, is the greatest common divisor of $\alpha, \beta, \ldots$}

2. \gesperrt{If $\alpha x + \beta y$ is a unit, then $\alpha$ and $\beta$ are called coprime or relatively prime. If there are functionals $\xi, \eta$ for which $\alpha \xi + \beta \eta$ is a unit, then $\alpha, \beta$ are relatively prime.}

3. \gesperrt{If $\alpha$ is relatively prime to $\beta$ and to $\gamma$, then it is also relatively prime to $\beta\gamma$.}

For by assumption
\[
\varepsilon = \alpha x + \beta y,\quad \varepsilon_1 = \alpha u + \beta v
% [sic] printed epsilon_1 = alpha u + beta v, should be alpha u + gamma v: alpha is assumed prime to gamma, and the next display expands (alpha x + beta y)(alpha u + gamma v)
\]
are units ($x, y, u, v$ new variables).
\[
\alpha(\alpha u x + \gamma v x + \beta u y) + \beta \gamma v y = \varepsilon \varepsilon_1,
\]
and hence, by 3., $\alpha$ and $\beta\gamma$ are relatively prime.
% [sic] printed "nach 3.", a circular reference to statement 3. itself, which is being proved; the step uses the criterion of 2. (alpha xi + beta eta a unit)

4. \gesperrt{If $\alpha$ is relatively prime to $\beta$ and $\alpha\mu$ is divisible by $\beta$, then $\mu$ is divisible by $\beta$.}

For from the assumption it follows that
\[
\alpha \mu x + \beta \mu y = \varepsilon \mu,
\]
from which, since $\varepsilon$ is a unit, the proof results.

\begin{center}
\S\,182. \textbf{Prime functionals in $\overline{\Omega}$.}
\end{center}

1. \gesperrt{Definition. An integral functional $\pi$ in $\overline{\Omega}$ is called a prime functional if it is not a unit and, apart from by the units, is divisible only by the functionals associated with it.}
